
\chapter{Buchstaben-Werkstatt II -- Die Sturköpfe}

Diese sechs Buchstaben sind Einzelgänger: \arbig{ا} \arbig{و} \arbig{ر}
\arbig{ز} \arbig{د} \arbig{ذ} geben zwar gerne eine Hand nach hinten,
weigern sich aber, nach vorne eine zweite Hand auszustrecken.

\begin{lachbox}
Das sind die Sturköpfe der Familie: sie hören zu (verbinden sich mit dem,
was davor kommt), aber sie antworten nie (verbinden sich nie mit dem, was
danach kommt). Ziemlich deutsches Verhalten, eigentlich.
\end{lachbox}

\section*{ا -- alif}
\begin{center}
\begin{tabular}{ll}
\textbf{Allein / Anfang:} & \arbig{ا} \\
\textbf{Nach einem Buchstaben:} & \textarabic{\Large بابا} \quad baaba (Papa) \\
\end{tabular}
\end{center}
\tracerow{ا}

\section*{و -- waaw}
\begin{center}
\begin{tabular}{ll}
\textbf{Allein / Anfang:} & \arbig{و} \quad \textarabic{\Large ورد} \quad ward (Rose) \\
\textbf{Nach einem Buchstaben:} & \textarabic{\Large أبو} \quad abu (Vater von...) \\
\end{tabular}
\end{center}
\tracerow{و}

\section*{ر -- raa}
\begin{center}
\begin{tabular}{ll}
\textbf{Allein / Anfang:} & \arbig{ر} \quad \textarabic{\Large ورد} \quad ward (Rose, r in der Mitte) \\
\textbf{Nach einem Buchstaben:} & \textarabic{\Large قمر} \quad qamar (Mond) \\
\end{tabular}
\end{center}
\tracerow{ر}

\section*{ز -- zaay}
\begin{center}
\begin{tabular}{ll}
\textbf{Allein / Anfang:} & \arbig{ز} \quad \textarabic{\Large زهرة} \quad zahra (Blume) \\
\textbf{Nach einem Buchstaben:} & \textarabic{\Large موز} \quad mooz (Banane) \\
\end{tabular}
\end{center}
\tracerow{ز}

\begin{tippbox}
موز (mooz) -- Banane. Ja, wirklich, es klingt fast wie im Deutschen mit
Akzent. Ägyptische Bananen sind kleiner und süßer -- probier sie.
\end{tippbox}

\section*{د -- daal}
\begin{center}
\begin{tabular}{ll}
\textbf{Allein / Anfang:} & \arbig{د} \quad \textarabic{\Large دار} \quad daar (Haus/Anwesen) \\
\textbf{Nach einem Buchstaben:} & \textarabic{\Large ولد} \quad walad (Junge) \\
\end{tabular}
\end{center}
\tracerow{د}

\section*{ذ -- dhaal}
Selten im Alltag -- im Masri wird daraus meist ein ,,d`` oder ,,z``.
\begin{center}
\begin{tabular}{ll}
\textbf{Allein / Anfang:} & \arbig{ذ} \\
\textbf{Nach einem Buchstaben:} & \textarabic{\Large هذا} \quad haadha (dieses hier; Masri oft: \emph{da}) \\
\end{tabular}
\end{center}
\tracerow{ذ}

\begin{checkpointbox}
Merksatz: \arbig{ا و ر ز د ذ} geben nie eine zweite Hand. Nach ihnen fängt
jeder Buchstabe wieder ganz von vorne an -- das ist Absicht, kein
Druckfehler!
\end{checkpointbox}
